\RequirePackage{iftex}
\ifPDFTeX
\documentclass[a4paper,12pt]{article}
\usepackage[margin=24mm]{geometry}
\usepackage[T1]{fontenc}
\usepackage[utf8]{inputenc}
\usepackage{CJKutf8}
\else
\documentclass[a4paper,12pt]{jsarticle}
\usepackage[dvipdfmx,margin=24mm]{geometry}
\fi
\usepackage{amsmath,amssymb}
\usepackage{enumitem}

\setlength{\parindent}{0pt}
\setlength{\parskip}{0.45em}
\setlist[enumerate]{label=(\alph*),leftmargin=3em,itemsep=0.2em,topsep=0.2em}

\newcommand{\ans}[1]{\par\vspace{0.8em}\noindent{\large 問題 #1}\quad}
\newcommand{\choiceans}[1]{\par\vspace{0.8em}\noindent{\large 選択問題 #1}\quad}

\begin{document}
\ifPDFTeX
\begin{CJK}{UTF8}{min}
\fi

\begin{center}
{\Large 数学1A 中間試験 解答}
\end{center}

\ans{1}
\begin{enumerate}
\item $\forall x\in\mathbb{R},\ x^2+1>0$.
\item $\exists n\in\mathbb{N}\ \text{such that}\ n^2>100$.
\end{enumerate}

\ans{2}
\begin{enumerate}
\item $x=1$ で連続であるためには
\[
\lim_{x\to 1-}f(x)=f(1)
\]
が必要である。左極限は $1+a$ であり, $f(1)=2$ だから
\[
1+a=2.
\]
よって $a=1$.

\item $f(x)=x^3+x-1$ は多項式なので $[0,1]$ で連続である。また
\[
f(0)=-1,\qquad f(1)=1.
\]
したがって中間値の定理より, ある $c\in(0,1)$ が存在して
\[
f(c)=0
\]
となる。よって $x^3+x-1=0$ は $(0,1)$ に少なくとも1つの解をもつ。
\end{enumerate}

\ans{3}
\begin{enumerate}
\item
\[
f'(x)=2x\sin x+(x^2+1)\cos x.
\]
\item
\[
g'(x)=\log x+1\qquad (x>0).
\]
\end{enumerate}

\ans{4}
分子, 分母はともに $x\to 0$ で $0$ に近づく。ロピタルの定理より
\[
\lim_{x\to 0}\frac{\cos(\sin x)-1}{x}
=\lim_{x\to 0}\frac{-\sin(\sin x)\cos x}{1}=0.
\]

\ans{5}
\[
\frac{\partial f}{\partial x}(x,y)=3x^2y^2+ye^{xy},
\qquad
\frac{\partial f}{\partial y}(x,y)=2x^3y+xe^{xy}.
\]

\ans{6}
$a=\frac{\pi}{2}$ とする。
\[
f(a)=1,\quad f'(a)=0,\quad f''(a)=-1,\quad f'''(a)=0.
\]
したがって3次のテイラー多項式は
\[
T_3(x)=1-\frac{1}{2}\left(x-\frac{\pi}{2}\right)^2.
\]

\ans{7}
\[
e^{2x}=\sum_{n=0}^{\infty}\frac{(2x)^n}{n!}
=1+2x+\frac{(2x)^2}{2!}+\frac{(2x)^3}{3!}+\cdots.
\]

\ans{8}
$T_n(x)=\sum_{k=0}^{n}\frac{x^k}{k!}$ である。
\begin{enumerate}
\item テイラーの定理より, ある $\theta\in(0,1)$ が存在して
\[
e=T_n(1)+\frac{e^\theta}{(n+1)!}.
\]
よって
\[
e-T_n(1)=\frac{e^\theta}{(n+1)!}>0.
\]
$e<3$ を用いると $e^\theta<e<3$ であり, $n>3$ ならば
\[
0<e-T_n(1)<\frac{3}{(n+1)!}<\frac{1}{n!}.
\]
したがって
\[
0<|e-T_n(1)|<\frac{1}{n!}.
\]

\item $e$ が有理数であると仮定し, $e=\frac{p}{q}$ と書く。ただし $p,q$ は自然数である。$n>3$ かつ $n\ge q$ となる $n$ をとる。

(a)より
\[
0<n!\left(e-T_n(1)\right)<1.
\]
一方,
\[
n!e=n!\frac{p}{q},\qquad
n!T_n(1)=\sum_{k=0}^{n}\frac{n!}{k!}
\]
はいずれも整数である。したがって $n!(e-T_n(1))$ は整数である。これは $0$ より大きく $1$ より小さい整数が存在することを意味し, 矛盾である。よって $e$ は無理数である。
\end{enumerate}

\ans{9}
\begin{enumerate}
\item オイラーの公式より
\[
\cos w=\frac{e^{iw}+e^{-iw}}{2}.
\]
$w=ix$ とすると
\[
\cos(ix)=\frac{e^{-x}+e^x}{2}.
\]

\item (a)より, $z=it$ とおくと
\[
\cos(it)=\frac{e^t+e^{-t}}{2}.
\]
これが $2$ となるには
\[
e^t+e^{-t}=4.
\]
$u=e^t$ とおくと
\[
u+\frac1u=4,\qquad u^2-4u+1=0.
\]
よって $u=2+\sqrt3$ とできるので
\[
z=i\log(2+\sqrt3)
\]
は $\cos z=2$ を満たす。

また, 一般に
\[
\cos z=2
\]
は
\[
\frac{e^{iz}+e^{-iz}}{2}=2
\]
と同値である。$u=e^{iz}$ とおくと
\[
u^2-4u+1=0
\]
となるので
\[
u=2\pm\sqrt3.
\]
したがって
\[
z=2\pi n\pm i\log(2+\sqrt3)\qquad (n\in\mathbb{Z})
\]
が $\cos z=2$ の複素解である。
\end{enumerate}

\ans{10}
\begin{enumerate}
\item
\[
f'(x)=\frac{\pi}{4}\cos\left(\frac{\pi}{4}x\right).
\]
$0<x<2$ では $0<\frac{\pi}{4}x<\frac{\pi}{2}$ だから
\[
\cos\left(\frac{\pi}{4}x\right)>0.
\]
よって $f'(x)>0$ である。したがって $f$ は $[0,2]$ 上で単調増加である。

\item まず $a_0=0\in[0,2]$ であり, $x\in[0,2]$ ならば
\[
1\le f(x)\le 2
\]
である。したがってすべての $n\ge 1$ について $a_n\in[0,2]$ である。

また $a_1=1\ge a_0$ であり, (a)より $f$ は単調増加だから, 帰納的に
\[
a_{n+1}=f(a_n)\ge f(a_{n-1})=a_n
\]
となる。よって $(a_n)$ は単調増加で, 上に $2$ で有界である。したがって極限 $L$ が存在する。

漸化式で極限をとると
\[
L=1+\sin\left(\frac{\pi}{4}L\right).
\]
$h(x)=1+\sin\left(\frac{\pi}{4}x\right)-x$ とおくと
\[
h'(x)=\frac{\pi}{4}\cos\left(\frac{\pi}{4}x\right)-1<0
\]
が $[0,2]$ で成り立つ。また $h(2)=0$ である。したがって $[0,2]$ における解は $L=2$ のみである。

よって
\[
\lim_{n\to\infty}a_n=2.
\]

別解. すでに示したように, $n\ge 1$ では $1\le a_n\le 2$ である。
$f(2)=2$ だから, 平均値の定理より, ある $\xi_n\in(a_n,2)$ が存在して
\[
2-a_{n+1}=f(2)-f(a_n)=f'(\xi_n)(2-a_n)
\]
となる。ここで $\xi_n\in[1,2]$ なので
\[
0\le f'(\xi_n)
=\frac{\pi}{4}\cos\left(\frac{\pi}{4}\xi_n\right)
\le \frac{\pi}{4}\cos\frac{\pi}{4}
=\frac{\pi\sqrt2}{8}<1.
\]
したがって
\[
0\le 2-a_{n+1}\le \frac{\pi\sqrt2}{8}(2-a_n)\qquad (n\ge 1)
\]
である。これを繰り返すと
\[
0\le 2-a_n\le
\left(\frac{\pi\sqrt2}{8}\right)^{n-1}(2-a_1)
\]
となる。右辺は $0$ に収束するので
\[
\lim_{n\to\infty}a_n=2.
\]
\end{enumerate}

\ifPDFTeX
\end{CJK}
\fi
\end{document}
